\section{Results}
\label{sec:resultados}

\Cref{tab:resultados} consolidates the performance of the three classifiers
on the 2025 validation together with the two baselines of the autoregressive
model, and \cref{tab:resultados-ejes} summarizes what each interdependence
technique delivered. The three classifiers exceed the 0.75 performance
threshold set in \cref{sec:obj:objetivo}.

\begin{table}[t]
  \centering
  \caption{Performance on the 2025 validation. The station-day models
    (logistic and discriminant) and the metropolitan-area model
    (autoregressive and its baselines) answer different questions, and their
    AUCs are not directly comparable across blocks. Sensitivity and
    specificity at the Youden threshold set on the fitting data;
    generalization is the out-of-sample AUC from leave-one-year-out
    cross-validation.}
  \label{tab:resultados}
  \small
  \begin{tabular}{lrrrr}
    \toprule
    Model & AUC & Sens. & Spec. & General. \\
    \midrule
    \multicolumn{5}{l}{\emph{Station-day} ($n=2{,}927$)} \\
    Logistic regression    & 0.850 & 0.698 & 0.853 & 0.80 \\
    Linear discriminant    & 0.846 & 0.758 & 0.788 & 0.80 \\
    \midrule
    \multicolumn{5}{l}{\emph{Metropolitan area} ($n=365$)} \\
    Autoregressive AR(1)   & 0.898 & 0.759 & 0.841 & 0.878 \\
    Climatology (baseline) & 0.668 & ---   & ---   & --- \\
    Persistence (baseline) & 0.717 & ---   & ---   & --- \\
    \bottomrule
  \end{tabular}
\end{table}

\begin{table}[t]
  \centering
  \caption{Axes retained by the interdependence techniques on the 2021--2024
    fitting data. Variance: total for PCA, common for FA.}
  \label{tab:resultados-ejes}
  \small
  \begin{tabular}{llr}
    \toprule
    Technique & Axis & Var. (\%) \\
    \midrule
    \multirow{3}{*}{PCA (15 variables)}
      & Primary combustion load      & 26.1 \\
      & Photochemical forcing        & 19.8 \\
      & Atmospheric stagnation       & 10.0 \\
    \midrule
    \multirow{3}{*}{FA (13 variables)}
      & Primary combustion           & 15.7 \\
      & Ventilation                  & 14.5 \\
      & Photochemical forcing        & 10.1 \\
    \bottomrule
  \end{tabular}
\end{table}

Four results cut across the techniques. First, the same three axes emerge
along two different paths, and FA corrects PCA on one point: the stagnation
axis was a wind axis, because barometric pressure has a communality of 0.13
and almost all of its variance is specific. Second, the hierarchy of effects
is stable across the three classifiers: photochemical forcing weighs more
than primary combustion, ventilation protects, and in the station-day models
season and zone weigh more than any factor. Third, temporal dependence is
real and exploitable: the response lag raises the metropolitan-area AUC from
0.882 to 0.898 and multiplies the odds by 3.83. Fourth, all models share a
year effect that no predictor explains: 2024 and 2025 are the highest-ozone
years of the series and also the best classified, and the performance
expected in an arbitrary year is between 0.02 and 0.05 below that of the
validation.
